\chapter[Performance Tuning Avanzado]{\emj{⚡}~Performance Tuning Avanzado}

\section[Autovacuum: el guardián silencioso]{\emj{🧹}~Autovacuum: el guardián silencioso}

MVCC (Multi-Version Concurrency Control) significa que cada UPDATE en PostgreSQL crea una nueva versión de la fila; la vieja queda como ``dead tuple''.
Autovacuum limpia dead tuples, actualiza estadísticas y previene \textbf{transaction ID wraparound} (el peor incidente de PostgreSQL: la BD se pone en modo solo-lectura).

\begin{teoriabox}{Cómo funciona autovacuum}
\begin{itemize}
  \item \textbf{autovacuum launcher:} proceso maestro, activa workers según tablas que lo necesitan
  \item \textbf{autovacuum worker:} trabaja una tabla a la vez; limitado por \texttt{autovacuum\_max\_workers} (default 3)
  \item \textbf{Trigger de vacuum:} tabla con $>$ \texttt{autovacuum\_vacuum\_threshold + autovacuum\_vacuum\_scale\_factor × n\_live\_tup} dead tuples
  \item \textbf{Trigger de analyze:} similar pero para estadísticas (\texttt{autovacuum\_analyze\_scale\_factor})
  \item Autovacuum usa \texttt{cost delay} para no saturar I/O: trabaja en ráfagas, luego duerme
\end{itemize}
\end{teoriabox}

\begin{lstlisting}[caption={Diagnosticar estado de autovacuum}]
-- Tablas que más necesitan vacuum:
SELECT schemaname, relname,
       n_dead_tup, n_live_tup,
       round(n_dead_tup::numeric / NULLIF(n_live_tup + n_dead_tup, 0) * 100, 2) AS dead_pct,
       last_autovacuum, last_autoanalyze,
       autovacuum_count, autoanalyze_count
FROM pg_stat_user_tables
ORDER BY n_dead_tup DESC
LIMIT 20;

-- Vacuums activos ahora mismo:
SELECT pid, query, state, now() - pg_stat_activity.query_start AS duration
FROM pg_stat_activity
WHERE query LIKE 'autovacuum%';

-- ¿Hay tablas con bloat alto?
SELECT relname, pg_size_pretty(pg_total_relation_size(oid)) AS total_size,
       pg_size_pretty(pg_relation_size(oid)) AS table_size
FROM pg_class WHERE relkind = 'r'
ORDER BY pg_total_relation_size(oid) DESC LIMIT 10;
\end{lstlisting}

\begin{lstlisting}[caption={Tuning de autovacuum por tabla}]
-- Para tablas con alta tasa de cambios (ej: logs, eventos):
ALTER TABLE events SET (
  autovacuum_vacuum_scale_factor = 0.01,   -- vacuum cuando 1% de filas son dead
  autovacuum_analyze_scale_factor = 0.005,
  autovacuum_vacuum_cost_delay = 2,        -- ms de delay entre ráfagas (default 20)
  autovacuum_vacuum_cost_limit = 400       -- trabajo por ráfaga (default 200)
);

-- VACUUM manual agresivo (para emergencias, bloquea concurrent vacuums):
VACUUM (ANALYZE, VERBOSE) orders;

-- VACUUM FULL: compacta tabla reescribiéndola (exclusivo, bloquea todo):
-- Solo en ventana de mantenimiento, o usar pg_repack
VACUUM FULL orders;
\end{lstlisting}

\section[Table Bloat y fill\_factor]{\emj{📦}~Table Bloat y fill\_factor}

Bloat ocurre cuando vacuum no puede reclamar espacio porque hay transacciones largas activas, o cuando el espacio muerto no se reutiliza por la estructura física de páginas.

El \texttt{fill\_factor} es la solución preventiva al bloat causado por UPDATEs frecuentes.
PostgreSQL organiza los datos en páginas de 8KB. Por defecto, cada página se llena al 100\% con filas nuevas.
Cuando se hace UPDATE, la fila vieja queda muerta en su página original, y la nueva se escribe en una página distinta si no hay espacio libre.
Esto obliga al índice a apuntar a la nueva ubicación (heap fetch costoso) y acumula dead tuples rápidamente.

Con \texttt{fill\_factor = 70}, PG deja 30\% de cada página libre al insertar.
Cuando llega un UPDATE, si la nueva versión de la fila cabe en el espacio libre de \emph{la misma página}, se realiza un \textbf{HOT update} (Heap-Only Tuple): no se actualiza el índice, no hay heap fetch adicional, y el dead tuple se puede reclamar en la misma página sin VACUUM completo.
El resultado: menos bloat, índices más estables, y updates más rápidos.

\begin{lstlisting}[caption={fill\_factor para tablas con muchos UPDATEs}]
-- fill_factor: qué porcentaje de cada página se llena al INSERT.
-- El resto se reserva para UPDATEs "in-page" (HOT updates).
-- HOT update: si la fila cabe en la misma página → no se actualiza el índice → mucho más rápido

-- Para tablas con muchos UPDATEs, reducir fill_factor:
ALTER TABLE orders SET (fillfactor = 70);
-- Necesita VACUUM o CLUSTER para aplicar en páginas existentes

-- Para tablas append-only (logs, eventos):
ALTER TABLE audit_log SET (fillfactor = 100);  -- default

-- Ver fill_factor actual:
SELECT relname, reloptions FROM pg_class
WHERE relname = 'orders';
\end{lstlisting}

\section[Connection Pooling con pgBouncer]{\emj{🔌}~Connection Pooling con pgBouncer}

Cada conexión de PostgreSQL es un proceso OS. Con 200+ conexiones directas, el overhead de context switching degrada el rendimiento. pgBouncer es un proxy que multiplexa muchas conexiones de cliente sobre pocas conexiones reales.

\begin{teoriabox}{Modos de pgBouncer}
\begin{itemize}
  \item \textbf{Session mode:} una conexión de PG por conexión de cliente mientras dure la sesión. Seguro pero no multiplexea mucho.
  \item \textbf{Transaction mode:} una conexión de PG por transacción. Multiplexeo real. \textbf{Limitación:} no se pueden usar prepared statements de sesión, \texttt{SET LOCAL}, \texttt{LISTEN/NOTIFY}.
  \item \textbf{Statement mode:} una conexión por statement. Solo para apps que no usan transacciones. Raramente útil.
\end{itemize}
\end{teoriabox}

\begin{lstlisting}[language=, caption={pgbouncer.ini — configuración típica}]
[databases]
mydb = host=127.0.0.1 port=5432 dbname=mydb

[pgbouncer]
listen_port = 6432
listen_addr = *
auth_type = scram-sha-256
auth_file = /etc/pgbouncer/userlist.txt

pool_mode = transaction
max_client_conn = 1000
default_pool_size = 20          ; conexiones reales a PG
min_pool_size = 5
reserve_pool_size = 5
reserve_pool_timeout = 3

server_idle_timeout = 600
client_idle_timeout = 0

; Stats cada 60s en pgbouncer log
stats_period = 60
\end{lstlisting}

\begin{lstlisting}[caption={Monitorear pgBouncer desde psql}]
-- Conectarse a la base admin de pgbouncer (puerto 6432):
-- psql -h localhost -p 6432 -U pgbouncer pgbouncer

SHOW POOLS;    -- ver cl_active, cl_waiting, sv_active, sv_idle por pool
SHOW STATS;    -- requests/s, bytes/s, query time promedio
SHOW CLIENTS;  -- conexiones de cliente activas
SHOW SERVERS;  -- conexiones a PostgreSQL activas
\end{lstlisting}

\section[Índices Avanzados]{\emj{📑}~Índices Avanzados}

Un índice estándar de B-tree en PostgreSQL indexa \emph{todas} las filas de una columna.
Pero existen tres variantes avanzadas que pueden ser órdenes de magnitud más eficientes dependiendo del patrón de acceso:

\begin{itemize}
  \item \textbf{Partial index:} solo indexa las filas que cumplen una condición WHERE. Si el 95\% de las órdenes están ``completed'' y solo consultas frecuentes son sobre ``pending'' (el 5\%), un índice parcial sobre pending es 20 veces más pequeño y se actualiza 20 veces menos.
  \item \textbf{Expression index:} indexa el resultado de una función, no el valor crudo. Útil cuando la query siempre aplica la misma transformación (ej: \texttt{lower(email)} para búsquedas case-insensitive). Sin este índice, la función se evalúa en cada fila durante el seq scan.
  \item \textbf{Covering index (INCLUDE):} incluye columnas adicionales dentro del índice sin que sean parte de la clave de búsqueda. Permite que el planner satisfaga la query completamente desde el índice (Index Only Scan), sin tocar el heap (la tabla principal). Esto elimina el costo de random I/O de los heap fetches.
\end{itemize}

\begin{lstlisting}[caption={Partial, expression y covering indexes}]
-- Partial index: solo indexa las filas que cumplen la condición
-- Mucho más pequeño y mantenido solo cuando aplica el WHERE
CREATE INDEX idx_orders_pending
ON orders (created_at)
WHERE status = 'pending';

-- Expression index: indexa el resultado de una expresión
CREATE INDEX idx_users_email_lower
ON users (lower(email));

-- Ahora solo usa el índice con: WHERE lower(email) = 'foo@example.com'
-- No con: WHERE email = 'foo@example.com' (string diferente)

-- Covering index (INCLUDE): añade columnas al índice sin que sean clave
-- Permite Index Only Scan sin ir al heap
CREATE INDEX idx_orders_customer_covering
ON orders (customer_id) INCLUDE (status, total_amount, created_at);

-- La query: SELECT status, total_amount FROM orders WHERE customer_id = 42
-- puede usar Index Only Scan si visibility map está al día

-- Concurrent index build (no bloquea escrituras):
CREATE INDEX CONCURRENTLY idx_items_order_id ON items (order_id);
-- Tarda más, pero la tabla sigue accesible
\end{lstlisting}

Un índice nunca usado no es gratuito: PostgreSQL debe actualizarlo en cada INSERT, UPDATE y DELETE sobre la tabla.
Un índice de 10GB que nadie usa consume esa RAM en el buffer cache, agrega write amplification en cada transacción, y ralentiza VACUUM.
Las siguientes queries identifican índices candidatos a eliminación y tablas que podrían beneficiarse de nuevos índices:

\begin{lstlisting}[caption={Detectar índices ineficientes}]
-- Índices nunca usados:
SELECT schemaname, relname, indexrelname, idx_scan, idx_tup_read
FROM pg_stat_user_indexes
WHERE idx_scan = 0
ORDER BY pg_relation_size(indexrelid) DESC;

-- Tablas con seq scans frecuentes (posibles índices faltantes):
SELECT relname, seq_scan, idx_scan,
       seq_scan::float / NULLIF(seq_scan + idx_scan, 0) AS seq_ratio
FROM pg_stat_user_tables
WHERE seq_scan > 100
ORDER BY seq_scan DESC;

-- Tamaño de índices:
SELECT relname AS table, indexrelname AS index,
       pg_size_pretty(pg_relation_size(indexrelid)) AS size
FROM pg_stat_user_indexes
ORDER BY pg_relation_size(indexrelid) DESC;
\end{lstlisting}

\section[work\_mem y shared\_buffers]{\emj{💾}~work\_mem y shared\_buffers}

La memoria es el parámetro de mayor impacto en el rendimiento de PostgreSQL.
A diferencia de otros sistemas, PG no tiene un único pool de memoria central: cada operación (sort, hash join, hash aggregate) toma su propia porción de \texttt{work\_mem}.
Esto significa que la memoria total consumida puede escalar con el número de conexiones y la complejidad de las queries.

Entender cada parámetro y su relación es fundamental antes de tocar \texttt{postgresql.conf}: un ajuste mal calculado de \texttt{work\_mem} puede llevar el servidor a OOM (Out of Memory) y matar el proceso de PG.

\begin{teoriabox}{Parámetros de memoria críticos}
\begin{itemize}
  \item \textbf{shared\_buffers:} caché de páginas de PG. Regla: 25--40\% de RAM. No poner $>$40\% porque el OS cache también ayuda y PG lo usa.
  \item \textbf{work\_mem:} RAM por operación de sort/hash. Multiplicado por número de sorts simultáneos × conexiones. Regla conservadora: RAM / max\_connections / 4.
  \item \textbf{maintenance\_work\_mem:} para VACUUM, CREATE INDEX, CLUSTER. Puede ser mayor (512MB--1GB en servidores grandes). No afecta conexiones normales.
  \item \textbf{wal\_buffers:} buffer para escritura de WAL. Default: 3\% de shared\_buffers (mínimo 32kB, máximo 64MB). Raramente necesita ajuste.
  \item \textbf{effective\_cache\_size:} no reserva RAM; solo le dice al planner cuánta RAM total (PG + OS) puede usarse como caché. Afecta decisiones de Index Scan. Poner 50--75\% de RAM total.
\end{itemize}
\end{teoriabox}

\begin{lstlisting}[caption={Configuración de memoria recomendada (servidor 32GB)}]
-- postgresql.conf:
shared_buffers = 8GB              -- 25% de 32GB
work_mem = 64MB                   -- conservador; subir por sesión si necesario
maintenance_work_mem = 1GB
effective_cache_size = 24GB       -- 75% de 32GB
wal_buffers = 64MB

-- Para una query pesada específica (analytics), subir work_mem temporalmente:
SET work_mem = '2GB';
SELECT /* heavy aggregation */ ...;
RESET work_mem;
\end{lstlisting}

\begin{entrevistabox}
\begin{enumerate}
  \item ¿Qué es table bloat y cómo lo causas involuntariamente?
  \item ¿Cuándo usarías pgBouncer en modo transaction vs session? ¿Qué pierdes en modo transaction?
  \item Explica qué es un HOT update y cómo fill\_factor lo facilita.
  \item ¿Por qué un índice nunca usado sigue costando dinero en producción?
  \item ¿Qué pasa si pones work\_mem = 2GB con 200 conexiones y cada una hace 3 sorts simultáneos?
  \item ¿Cómo detectas si autovacuum no está manteniendo una tabla?
\end{enumerate}
\end{entrevistabox}
